% A proof that the genus-two star X_0^2002(1)^* is the curve C stored in data/models/models_2002_1.m,
% and what it gives for the genus-four quotient by <w_2, w_13, w_77> (Section 3).
%
% Every computation the proof quotes is a script in tools/proof-2002/, each printing the number
% used here.  Two of them need packages outside this repository (Genus2Conductor and
% endomorphisms, bundled in edgarcosta/CHIMP); the README there says which.
%
% Numbering is local to this file.
\documentclass[11pt,reqno]{amsart}
\usepackage[margin=1.15in]{geometry}
\usepackage{amsmath,amssymb,amsthm,mathtools}
\usepackage{booktabs}
\usepackage[colorlinks=true,linkcolor=blue,citecolor=blue,urlcolor=blue]{hyperref}
\theoremstyle{plain}
\newtheorem{theorem}{Theorem}
\newtheorem{proposition}[theorem]{Proposition}
\newtheorem{lemma}[theorem]{Lemma}
\theoremstyle{definition}
\newtheorem{remark}[theorem]{Remark}
\DeclareMathOperator{\Jac}{Jac}
\DeclareMathOperator{\End}{End}
\DeclareMathOperator{\Gal}{Gal}
\DeclareMathOperator{\Tr}{Tr}
\DeclareMathOperator{\Norm}{N}
\DeclareMathOperator{\Frob}{Frob}
\newcommand{\Q}{\mathbb{Q}}
\newcommand{\Z}{\mathbb{Z}}
\newcommand{\F}{\mathbb{F}}
\newcommand{\Qbar}{\overline{\mathbb{Q}}}
\newcommand{\OK}{\mathcal{O}_K}
\begin{document}
\title{The genus-two Shimura curve quotient $X_0^{2002}(1)^*$}
\author{}
\date{\today}
\maketitle

\section{The statement}

Let $X = X_0^{2002}(1)^*$ be the quotient of the Shimura curve of discriminant $2002 = 2\cdot 7\cdot
11\cdot 13$ and level one by its full group of Atkin--Lehner involutions. It has genus two. Let
\[
  C:\quad y^2 + (x^2+x)\,y \;=\; 16x^6 + 4x^5 + 180x^4 - 456x^3 + 196x^2 + 64x + 4 .
\]

\begin{theorem}\label{thm:main}
$X \cong C$ over $\Q$.
\end{theorem}

The curve $C$ was found twice, by two numerical routes sharing nothing: from the periods of the
modular abelian surface attached to the newform below, and from the differentials of $X$ evaluated
by the Shimizu lift. Neither route proves anything, and this note does not use how $C$ was found.
What it uses is that $C$ is in hand: every step below is either a theorem from the literature or a
finite exact computation on $C$ or on a space of modular symbols, and the computations are listed
in Section~\ref{sec:computations}.

Throughout, $f$ is the newform orbit \texttt{2002.2.a.f} of weight two, level $2002$ and trivial
character: its coefficient field is $K=\Q(\sqrt5)$ with ring of integers $\OK=\Z[\beta]$,
$\beta=(1+\sqrt5)/2$, it has no complex multiplication and no inner twist, and its Atkin--Lehner
sign is $-1$ at each of $2,7,11,13$ \cite{lmfdb}. We write $a_p\in\OK$ for its coefficients,
$\sigma$ for the nontrivial automorphism of $K$, $t_p=\Tr_{K/\Q}a_p$ and $n_p=\Norm_{K/\Q}a_p$,
and $A_f$ for the abelian surface over $\Q$ attached to $f$ by Shimura.

\section{The proof}

\subsection*{Step 1: $\Jac(X)$ is $\Q$-isogenous to $A_f$.}
By Jacquet--Langlands the weight-two forms on $X_0^{2002}(1)$ correspond to the newforms of level
$2002$, with the Atkin--Lehner eigenvalue at each prime dividing the discriminant reversed; the
differentials on the quotient $X$ are those fixed by every Atkin--Lehner involution, hence
correspond to the newforms whose classical sign is $-1$ at each of $2,7,11,13$. A search over the
newforms of level $2002$ shows this picks out the single orbit $f$, of dimension two. By
Eichler--Shimura for Shimura curves \cite{ribet80}, $\Jac(X)\sim_\Q A_f$.

\subsection*{Step 2: $\End_\Q(\Jac C)=\OK$.}
The algorithm of Costa, Mascot, Sijsling and Voight \cite{cmsv} computes the endomorphism ring of a
Jacobian numerically and then \emph{certifies} the answer: it verifies that the lattice found is
saturated, and it realises each generator as an explicit correspondence on the curve by exact
computation. On $C$ it returns the generators $1$ and an endomorphism whose action on the
differentials has matrix $\begin{psmallmatrix}0&1\\1&-1\end{psmallmatrix}$, of characteristic
polynomial $x^2+x-1$, so that the ring generated is $\Z[(-1+\sqrt5)/2]=\OK$; the correspondence
realising it is written out in the certificate. So $\Jac(C)$ is of $\mathrm{GL}_2$-type over $\Q$
with $\End_\Q(\Jac C)=\OK$, and the same holds over $\Qbar$.

\subsection*{Step 3: the conductor of $\Jac(C)$ is $2002^2$.}
The discriminant of $C$ is $2^{20}\,7^6\,11^4\,13^6$, so $\Jac(C)$ has good reduction outside
$\{2,7,11,13\}$. The conductor exponents at $7$, $11$, $13$ are $2$ each, by Liu's algorithm as
implemented in Magma and in \textsc{Pari}'s \texttt{genus2red}, which agree. At $2$ the usual
formula of Ogg is not a theorem in genus two --- Liu gave counterexamples --- and Magma says so when
it falls back on it; the exponent at $2$ is instead computed by the $3$-torsion method of Dokchitser
and Doris \cite{dd}, with its proof flag set, and is $2$. Hence the conductor is
$2^2 7^2 11^2 13^2 = 2002^2$.

\subsection*{Step 4: $\Jac(C)$ is $\Q$-isogenous to $A_f$.}
By Ribet \cite[Theorem~4.4]{ribet92}, now unconditional through the proof of Serre's conjecture by
Khare and Wintenberger \cite{kw}, an abelian variety of $\mathrm{GL}_2$-type over $\Q$ is
$\Q$-isogenous to $A_g$ for a newform $g$ of weight two, whose level $N$ satisfies
$\mathrm{cond}(A_g)=N^{\dim A_g}$ \cite{carayol}. By Steps 2 and 3, $g$ has level $2002$ and
coefficient field $\Q(\sqrt5)=\End_\Q(\Jac C)\otimes\Q$; a newform whose coefficient field is
totally real has trivial character. The Euler factor of $\Jac(C)$ at a good prime $p$ is that of
$A_g$, namely $\prod_{\tau}(1-\tau(a_p(g))T+pT^2)$ over the two embeddings of $\Q(\sqrt5)$. There
are five Galois orbits of newforms of level $2002$, weight two and trivial character with
coefficient field of degree two (three with field $\Q(\sqrt5)$, two with $\Q(\sqrt{13})$), and
exactly one of them has the Euler factor of $\Jac(C)$ at $p=3$, namely
$1+3T+7T^2+9T^3+9T^4$; it is $f$, and the other four differ already there. So $g=f$ up to
conjugation, and $\Jac(C)\sim_\Q A_f$. (The Euler factors agree at every good prime below $120$;
one prime is what the argument needs.)

\subsection*{Step 5: every $\Q$-isogeny out of $\Jac(C)$ is an isomorphism.}
Write $J=\Jac(C)$ and, for a prime $\lambda$ of $\OK$, let $\rho_\lambda$ be the representation of
$\Gal(\Qbar/\Q)$ on $J[\lambda]$, a two-dimensional space over $\F_\lambda=\OK/\lambda$; its
semisimplification is an isogeny invariant, so it is also that of $A_f[\lambda]$, whose Frobenius
traces are $a_p \bmod \lambda$.

\begin{lemma}\label{lem:irred}
$\rho_\lambda$ is irreducible for every prime $\lambda$ of $\OK$.
\end{lemma}
\begin{proof}
Suppose $\rho_\lambda$ reducible, $\lambda\mid\ell$. Its semisimplification is $\chi_1\oplus\chi_2$
with $\chi_1\chi_2$ the mod-$\ell$ cyclotomic character, both unramified outside $2002\ell$, and
both of order prime to $\ell$ since their values lie in $\F_\lambda^\times$; so neither is wildly
ramified at $\ell$, and each is $\varepsilon_i\,\omega^{a_i}$ with $\omega$ the cyclotomic
character and $\varepsilon_i$ of conductor prime to $\ell$. The prime-to-$\ell$ Artin
conductor of $\rho_\lambda$ divides that of the $\lambda$-adic representation \cite{carayol89},
which is the prime-to-$\ell$ part of $2002$; for a reducible $\rho_\lambda$ that conductor is
$\mathrm{cond}(\varepsilon_1)\mathrm{cond}(\varepsilon_2)$, so each $\varepsilon_i$ has conductor
dividing $2002$, hence dividing $1001$, there being no nontrivial character of conductor $2$. In
particular $\varepsilon_i(p)=1$ at every prime $p\equiv1\pmod{2002}$. It remains to pin the
exponents $a_i$, and there are three cases.

\emph{(i) $\ell\nmid2002$, which includes $\ell=5$.} In weight two with $\ell$ prime to the level,
the restriction of $\rho_\lambda$ to the inertia group at $\ell$ is either absolutely irreducible,
given by the two fundamental characters of level two, or has semisimplification $1\oplus\omega$
\cite{edixhoven}. A globally reducible $\rho_\lambda$ is in the second case, so
$\{a_1,a_2\}=\{0,1\}$, and at $p\equiv1\pmod{2002}$ we get $a_p\equiv p^{a_1}+p^{a_2}=1+p
\pmod\lambda$.

\emph{(ii) $\ell\in\{7,11,13\}$.} Here $\ell$ divides the level exactly once and the character is
trivial, so the $\lambda$-adic representation is special (Steinberg) at $\ell$: its restriction to
the decomposition group is an unramified twist of $\begin{psmallmatrix}\tilde\omega&*\\0&1
\end{psmallmatrix}$, the twist being the Atkin--Lehner sign. Reducing, the restriction of
$\rho_\lambda$ to inertia at $\ell$ again has semisimplification $\omega\oplus1$, so
$\{a_1,a_2\}=\{0,1\}$ and the conclusion is the same as in (i).

\emph{(iii) $\ell=2$.} The prime $2$ is inert in $K$, so $\lambda=(2)$ and $\F_\lambda=\F_4$, whose
multiplicative group has order $3$; thus $\chi_1,\chi_2$ are Dirichlet characters of order dividing
$3$. Such a character is unramified at $2$, since $(\Z/2^k)^\times$ is a $2$-group for every $k$,
and $\omega\equiv1\pmod 2$; so $\chi_i=\varepsilon_i$ with conductor dividing $1001$ as above, and at
$p\equiv1\pmod{2002}$ we get $a_p\equiv\chi_1(p)+\chi_2(p)=2\equiv0\equiv1+p\pmod\lambda$, $p$
being odd. The congruence is the same one.

In every case, then, $\lambda$ divides $\Norm_{K/\Q}(a_p-1-p)$ for every prime
$p\equiv1\pmod{2002}$. Over the $24$ primes $p\equiv1\pmod{2002}$
below $200000$ the greatest common divisor of these norms is $1$.
\end{proof}

\begin{lemma}\label{lem:conj}
For every rational prime $\ell$, the $\F_\ell[\Gal(\Qbar/\Q)]$-module $J[\ell]$ is a direct sum of
pairwise non-isomorphic irreducible modules, except that $J[5]$ is a non-split extension of
$\rho_{(\sqrt5)}$ by itself.
\end{lemma}
\begin{proof}
If $\ell$ splits, $J[\ell]=\rho_\lambda\oplus\rho_{\lambda'}$ with $\F_\lambda=\F_\ell$, each
irreducible by Lemma~\ref{lem:irred}, and $\rho_\lambda\cong\rho_{\lambda'}$ would force $a_p\equiv
\sigma(a_p)\pmod\lambda$ for all $p$. If $\ell$ is inert, $J[\ell]=\rho_\lambda$ is a
two-dimensional $\F_{\ell^2}$-space, and the restriction of scalars of an irreducible
$\F_{\ell^2}$-representation to $\F_\ell$ is irreducible exactly when the representation is not
isomorphic to its conjugate under $\Gal(\F_{\ell^2}/\F_\ell)$, which acts on $\OK/\ell$ as
$\sigma$; again that isomorphism would force $a_p\equiv\sigma(a_p)\pmod\lambda$. In either case
$\lambda$ would divide $(a_p-\sigma(a_p))^2$ for every good $p$. Over the good primes below $400$
the greatest common divisor of $(a_p-\sigma(a_p))^2 = t_p^2-4n_p$ is exactly $5$ --- the factor
forced by $a_p-\sigma(a_p)\in\sqrt5\,\OK$ --- so no such $\lambda$ exists.

At $\ell=5$, $(\sqrt5)^2=(5)$ and $J[5]$ is an extension of $\rho=\rho_{(\sqrt5)}$ by itself. Were
it split as a Galois module, $J[5]\cong\rho\oplus\rho$, then for every good prime $p$ the
$\Frob_p$-invariants $J(\F_p)[5]$ would have even $\F_5$-dimension. But $J(\F_{29})\cong
\Z/4\times\Z/220$, of $5$-rank one. (The same happens at $37, 71, 101, 103, 113, 149$.)
\end{proof}

\begin{proposition}\label{prop:iso}
Every finite subgroup of $J(\Qbar)$ stable under $\Gal(\Qbar/\Q)$ is $J[\mathfrak a]$ for an ideal
$\mathfrak a$ of $\OK$. Consequently every $\Q$-isogeny $J\to B$ is an isomorphism.
\end{proposition}
\begin{proof}
Let $G$ be such a subgroup. It is the direct sum of its $\ell$-primary parts, each Galois-stable, so
we may take $G\subset J[\ell^\infty]$ and argue by induction on $|G|$, the case $G=0$ being
trivial. If $G\neq0$ then $G\cap J[\ell]$ is a nonzero Galois-stable subgroup of $J[\ell]$. For
$\ell\ne5$, Lemma~\ref{lem:conj} and the fact that a submodule of a direct sum of pairwise
non-isomorphic irreducibles is a sum of some of them give $G\cap J[\ell]=J[\mathfrak a_1]$ for an
ideal $\mathfrak a_1\supsetneq(\ell)$ dividing $\ell$; for $\ell=5$ the Galois-stable subgroups of a
non-split self-extension of an irreducible are $0$, $J[\sqrt5]$ and $J[5]$, so the same holds. As
$K$ has class number one, $\mathfrak a_1=(\alpha_1)$. The endomorphism $\alpha_1$ has kernel
$J[\mathfrak a_1]\subset G$, so $\alpha_1(G)$ is a Galois-stable subgroup of $J[\ell^\infty]$ of
order $|G|/|J[\mathfrak a_1]|<|G|$; by induction $\alpha_1(G)=J[\mathfrak a']$, and then
$G=\{x: \alpha_1x\in J[\mathfrak a']\}=J[\mathfrak a_1\mathfrak a']$.

For the last sentence, the kernel of a $\Q$-isogeny $J\to B$ is $J[\mathfrak a]$ with $\mathfrak
a=(\alpha)$, and $\alpha$ induces an isomorphism $J/J[\mathfrak a]\to J$; so $B\cong J$.
\end{proof}

\subsection*{Step 6: $\Jac(X)\cong\Jac(C)$ over $\Q$ as unpolarised surfaces.}
By Steps 1 and 4 there is a $\Q$-isogeny $\psi\colon\Jac(X)\to\Jac(C)$. Both are principally
polarised, so the dual isogeny is a $\Q$-isogeny $\Jac(C)\to\Jac(X)$, which by
Proposition~\ref{prop:iso} is an isomorphism.

\subsection*{Step 7: the polarisation, and the curve.}
The surface $A=\Jac(C)$ has $\End_\Q A=\OK$, an order in a real quadratic field fixed pointwise by
complex conjugation, and $\OK^\times$ contains the unit $\beta$ of norm $-1$. By Gonz\'alez,
Gu\`ardia and Rotger \cite[Theorem~1.1(1)]{ggr}, \emph{a principally polarisable abelian surface
over $\Q$ with $\End_\Q A=R$ an order in a quadratic field has a single isomorphism class of
principal polarisations over $\Q$ whenever $R^\times$ contains a unit of negative norm}. The
canonical polarisations of $\Jac(X)$ and of $\Jac(C)$ are both principal polarisations of $A$
defined over $\Q$, so they are $\Q$-isomorphic: $(\Jac X,\Theta_X)\cong(\Jac C,\Theta_C)$ over
$\Q$. Finally, Torelli holds over the base field for hyperelliptic curves \cite{serre-lauter}:
since $-1$ is induced by the hyperelliptic involution, an isomorphism of principally polarised
Jacobians over $\Q$ comes from an isomorphism of the curves over $\Q$. Hence $X\cong C$ over $\Q$.
\qed

\begin{remark}
Step 7 is where the method of Section~1 of the period note would have failed by itself. Point
counts cannot distinguish curves whose Jacobians are isomorphic as unpolarised surfaces, and for a
field whose fundamental unit has norm $+1$ the theorem cited gives \emph{two} classes of principal
polarisations, hence possibly two non-isomorphic curves with the same counts at every prime ---
this is what happens at level $35$. For $\Q(\sqrt5)$ the fundamental unit has norm $-1$ and the
ambiguity is absent. Nothing in the proof uses the quadratic twist: Step 7 identifies $X$ with $C$
over $\Q$ outright, which agrees with the twisting constant $1$ that the counts return.
\end{remark}

\begin{remark}
The proof does not use Faltings--Serre. For a surface with $\End=\Z$ that method would be the way
to Step 4, and it is implemented in \cite{bppt}; but it requires the residual representation into
$\mathrm{GSp}_4(\F_2)$ to be absolutely irreducible, which fails for a $\mathrm{GL}_2$-type surface,
whose $2$-torsion splits over $\F_4$ into two pieces. For $\mathrm{GL}_2$-type the route through
Serre's conjecture is both shorter and standard, and its one non-formal input, the endomorphism
ring, is exactly what \cite{cmsv} certifies.
\end{remark}

\section{The genus-four quotient by $\langle w_2,w_{13},w_{77}\rangle$}\label{sec:genus4}

Let $G\cong(\Z/2)^4$ be the Atkin--Lehner group of $X_0^{2002}(1)$, let
$H=\langle w_2,w_{13},w_{77}\rangle=\{w_m : m\in\{1,2,13,26,77,154,1001,2002\}\}$, of index two in
$G$, and let $Y=X_0^{2002}(1)/H$. It has genus four and is one of the quotients left open by the
classification; Noether's count on its Shimizu-lifted differentials says, numerically, that it is
hyperelliptic. The involution $\bar w$ of $Y$ induced by $G/H$ has quotient $Y/\bar w=X\cong C$, so
$Y$ is a double cover of the curve of Theorem~\ref{thm:main}, and this section shows how far that
theorem carries. Put
\begin{align*}
  Y_2:&\quad y^2=-1001x^{10}+98289x^8-279360x^6+185856x^4-16384x^2+65536,\\
  C':&\quad y^2+(x^2+x)\,y=-2x^5+10x^4-44x^3+29x^2+9x+1 .
\end{align*}
$Y_2$ is hyperelliptic of genus four; its involutions $x\mapsto-x$ and $(x,y)\mapsto(-x,-y)$ have
quotients $\Q$-isomorphic to $C$ and to $C'$ respectively, so $\Jac(Y_2)\sim\Jac(C)\times\Jac(C')$,
and $C'$ has conductor $2002^2$.

\begin{theorem}\label{thm:genus4}
Let $\pi\colon Y\to C$ be the quotient by $\bar w$ and $B$ its branch divisor.
\begin{enumerate}
\item $B=P_{11}+P_{91}$, where $P_{11}$ and $P_{91}$ are the images in $C$ of the CM points of
  discriminants $-11$ and $-91$. Both are $\Q$-rational points of $C$.
\item $Y$ is hyperelliptic if and only if $P_{91}=\bar\iota(P_{11})$, $\bar\iota$ the hyperelliptic
  involution of $C$; that is, if and only if $B$ is a fibre of the hyperelliptic map of $C$.
\item If $B$ is a fibre of the hyperelliptic map of $C$, then it is the fibre at infinity, the two
  points $(1:\pm4:0)$, and $Y\cong Y_2$ over $\Q$.
\end{enumerate}
In particular: if $Y$ is hyperelliptic, then $Y\cong Y_2$ over $\Q$.
\end{theorem}

Everything in the theorem is exact. What it does not decide is whether $Y$ is hyperelliptic at
all, and by (2) that is the position of one rational divisor on $C$: whether $P_{11}+P_{91}$ is
linearly equivalent to $2W$ for a Weierstrass point $W$. Proposition~\ref{prop:count} below reduces
it to the existence of one Borcherds form on $C$, whose search is in progress; the numerical
evidence is in Remark~\ref{rem:gap}.

\subsection*{The branch locus}
A point of $C=X_0^{2002}(1)/G$ is a branch point of $\pi$ exactly when its $G$-orbit upstairs is a
single $H$-orbit, that is, when the stabiliser in $G$ of a point of that orbit is not contained in
$H$: when the orbit consists of fixed points of some $w_m$ with $m\notin H$. By Ogg
\cite{ogg}, the fixed points of $w_m$ on $X_0^D(1)$ are the CM points by the quadratic orders of
discriminant $-4m$ and, when $m\equiv3\pmod4$, $-m$, and the order $R$ contributes
$h(R)\prod_{p\mid D,\,p\nmid m}\bigl(1-\{R/p\}\bigr)$ of them, $\{R/p\}$ the Eichler symbol. For
the eight $m\notin H$ this gives (\texttt{stepH-branch.m}): $w_7$, $w_{14}$, $w_{22}$, $w_{143}$,
$w_{182}$ and $w_{286}$ have no fixed points, at each a prime of $2002$ splitting in the order;
$w_{11}$ has exactly eight, all of discriminant $-11$ (the order of discriminant $-44$ has conductor
$2$, and $2$ ramifies in the quaternion algebra); and $w_{91}$ has exactly eight, all of discriminant
$-91$. So the branch points of $\pi$ are the images of these sixteen points. Riemann--Hurwitz for
$\pi$ reads $2\cdot4-2=2(2\cdot2-2)+\deg B$, so $\deg B=2$; each discriminant contributes at least one
branch point, hence exactly one: the eight points of discriminant $-11$ form a single $G$-orbit and
so do the eight of discriminant $-91$. The set of CM points of a given discriminant is stable under
$\Gal(\Qbar/\Q)$, so each orbit is Galois-stable and its image is a rational point of $C$. This is
(1).

\subsection*{Hyperellipticity, and the normal form}
\begin{lemma}\label{lem:normalform}
Let $\pi\colon Y\to C$ be a double cover over $\Q$ of a curve of genus two by a curve of genus
four, with covering involution $\bar w$ and branch divisor $B=P+Q$, $P,Q\in C(\Q)$. If $Y$ is
hyperelliptic then $Q=\bar\iota(P)$, $C$ has a $\Q$-rational Weierstrass point, and there are a
coordinate $t$ on $Y/\iota=\mathbb P^1$ and a quintic $h_1\in\Q[u]$ with $h_1(0)\ne0$ such that
\[
  Y:\ y^2=h_1(t^2),\qquad \bar w\colon t\mapsto -t,\qquad C=Y/\bar w:\ y^2=h_1(u),\ u=t^2,
\]
the branch points being the two points of $C$ over $u=0$ and the rational Weierstrass point the
one over $u=\infty$. Conversely, if $C$ has a model $y^2=h_1(u)$ with $h_1$ a quintic, $h_1(0)\neq
0$, then for every $c\in\Q^\times$ the curve $Y_c\colon y^2=h_1(t^2/c)$ is a hyperelliptic double
cover of $C$, branched along the fibre $u=0$, and its second involution $(t,y)\mapsto(-t,-y)$ has
quotient $v^2=c\,u\,h_1(u)$, $v=ty$.
\end{lemma}

\begin{proof}
Let $\iota$ be the hyperelliptic involution of $Y$. It is central in $\operatorname{Aut}(Y)$, so it
commutes with $\bar w$ and descends to an involution of $C$ whose quotient is
$Y/\langle\iota,\bar w\rangle=(Y/\iota)/\bar w$, of genus zero; so it is $\bar\iota$. Let $\tilde
P,\tilde Q\in Y$ be the fixed points of $\bar w$, the points over $P$ and $Q$; $\iota$ permutes them.
The involution that $\bar w$ induces on $Y/\iota=\mathbb P^1$ is not the identity (otherwise $\bar w
\in\{1,\iota\}$, whereas $Y/\bar w$ has genus two), so it has exactly two fixed points $a,b\in
\mathbb P^1$, and $\tilde P,\tilde Q$ lie over them. Neither $a$ nor $b$ is a branch point of
$Y\to\mathbb P^1$: if, say, $a$ were, choose the coordinate $t$ with $a=0$, $b=\infty$ and write
$Y\colon y^2=F(t)$ with $F$ squarefree and $F(0)=0$; the lift of $t\mapsto-t$ to $Y$ is
$y\mapsto\lambda(t)y$ with $\lambda^2=F(-t)/F(t)$, but $F(-t)/F(t)$ has odd order at $t=0$ and is
not a square. So the fibres over $a$ and $b$ have two points each, $\bar w$ acts on each fibre, and
since it has exactly two fixed points on $Y$ it fixes one fibre pointwise and swaps the other.
Hence $\{\tilde P,\tilde Q\}$ is a fibre of $Y\to\mathbb P^1$, $\tilde Q=\iota\tilde P$ and
$Q=\bar\iota P$.

Now take $t$ with $t(\tilde P)=0$, which is a rational coordinate because $\tilde P$ is a
rational point, and $b=\infty$ (rational, being the other fixed point of a rational involution of
$\mathbb P^1$ with one rational fixed point). Then $\bar w$ is $t\mapsto-t$, $y\mapsto\pm y$, and
since $\bar w$ fixes the points over $t=0$ the sign is $+$; so $F(-t)=F(t)$, $F$ is even, and being
squarefree of degree $9$ or $10$ it has degree $10$: $F(t)=h_1(t^2)$ with $\deg h_1=5$ and
$h_1(0)=F(0)\neq0$. The invariants of $\bar w$ are $u=t^2$ and $y$, so $C=Y/\bar w$ is $y^2=h_1(u)$,
with the branch points over $u=0$ and a rational Weierstrass point over $u=\infty$, $h_1$ having
odd degree. The converse is a computation: $Y_c$ has genus four, $t\mapsto-t$ fixes exactly the
two points over $t=0$, its quotient is $y^2=h_1(u)$ with $u=t^2/c$, and the invariants of
$(t,y)\mapsto(-t,-y)$ are $u$ and $v=ty$, with $v^2=t^2h_1(t^2/c)=c\,u\,h_1(u)$.
\end{proof}

Part (2) of the theorem is the first assertion of the lemma together with its converse: if $P_{91}=
\bar\iota P_{11}$ then $B$ is a fibre $\{u=u_0\}$ of the hyperelliptic map, and $C$ has a rational
Weierstrass point (below), so after moving that point to $u=\infty$ and $u_0$ to $0$ the lemma's
converse applies to some double cover branched along $B$; that $Y$ itself is such a cover, hence
hyperelliptic, is the argument of the next paragraph, which shows that all of them are.

\subsection*{The covers branched along a fibre}
The sextic $4f+h^2=64x^6+16x^5+721x^4-1822x^3+785x^2+256x+16$ of $C$ factors over $\Q$ as
$(8x+1)\,q(x)$ with $q$ irreducible of degree five, so $C$ has exactly one rational Weierstrass
point, $W$, over $x=-1/8$. The $2$-torsion of $\Jac(C)$ is spanned by the differences of
Weierstrass points, and a rational $2$-torsion class corresponds to a Galois-stable even subset of
the six roots up to complement; with Galois orbits of sizes $1$ and $5$ there is none, so
$\Jac(C)(\Q)[2]=0$.

Suppose $B$ is a fibre of the hyperelliptic map, over the point $x_0\in\mathbb P^1(\Q)$ (we
write $g=4f+h^2$ and work on $y^2=g(x)$). Put $u=(x-x_0)/(x+\tfrac18)$, or $u=8/(8x+1)$ if
$x_0=\infty$, so that $u(W)=\infty$ and $u=0$ along $B$; then $C$ is $y^2=q_1(u)$ with $q_1$ the
quintic $(u-1)^6g\bigl((-\tfrac18u-x_0)/(u-1)\bigr)$ up to a square constant, the degree dropping
because $W$ has gone to infinity. For $x_0=\infty$,
\[
  q_1(u)=64\,u^6\,g(-\tfrac18+\tfrac1u)=-2002u^5+98289u^4-139680u^3+46464u^2-2048u+4096 .
\]
Any double cover $Y'\to C$ over $\Q$ has function field $\Q(C)(\sqrt{\gamma})$ for some
$\gamma\in\Q(C)^\times$, and its branch divisor is the odd part of $\operatorname{div}\gamma$. The
function $u$ has divisor $B-2W$. If $Y'$ is branched along $B$ then $\operatorname{div}(\gamma/u)$ is
even, say $2D$; the class of $D$ lies in $\Jac(C)(\Q)[2]=0$, so $D=\operatorname{div}\delta$ with
$\delta\in\Q(C)^\times$ and $\gamma/u\delta^2$ is a constant $c\in\Q^\times$. Hence every double
cover of $C$ over $\Q$ branched along $B$ is one of the $Y_c\colon y^2=q_1(t^2/c)$ of
Lemma~\ref{lem:normalform}, $c\in\Q^\times/\Q^{\times2}$, and $Y$ is one of them. Its $c$ is
supported on $\{-1,2,7,11,13\}$: $Y$ has good reduction at every $p\nmid2002$ (Morita's theorem for
$X_0^{2002}(1)$, and the quotient by a group of order prime to $p$ of a curve with good reduction has
good reduction), so $\Jac(Y)\sim\Jac(C)\times\Jac(C'_c)$ does, where $C'_c\colon v^2=c\,u\,q_1(u)$
is the quadratic twist by $c$ of $C'_1$; a twist by a character ramified at an odd prime $p$ of good
reduction has bad reduction at $p$, so no such $p$ divides $c$ to an odd power. That leaves $32$
curves for each fibre, and the trace formula gives $\#Y(\F_p)$ at every good prime, as $p+1$ minus
the trace of Frobenius on the $H$-invariants of the weight-two forms of level $2002$ (the same
computation that underlies Step~1). By the Klein-four structure $\#Y_c(\F_p)=\#C(\F_p)+\#C'_c(\F_p)
-(p+1)$, so the candidates are told apart by the genus-two curves $C'_c$ alone.

\subsection*{The branch fibre}
It remains to find the fibre, and $C(\Q)$ cannot simply be listed: $\Jac(C)$ has rank $2$, equal to
the genus, so Chabauty's method does not apply and $C(\Q)$ is infinite. What makes the search
finite is reduction.

\begin{proposition}\label{prop:fibre}
Let $S=\{2,7,11,13\}$. If $B$ is a fibre $\{(a:\pm y':b)\}$ of $C$, written in the weighted
coordinates of $y^2=g(x)$ with $\gcd(a,b)=1$, then $y'$ is an $S$-unit. There are exactly seven such
fibres: $x=\infty,-1,0,\tfrac32,\tfrac13,-\tfrac14,\tfrac54$.
\end{proposition}

\begin{proof}
Let $p\notin S$. Then $Y$ and $C$ have good reduction at $p$ and $\pi$ extends to the smooth
models, as the quotient by $\bar w$; since $p$ is odd the reduction $\pi_p\colon Y_p\to C_p$ is a
tame double cover of smooth curves of genera $4$ and $2$, so by Riemann--Hurwitz it has exactly two
ramification points, and they lie in distinct fibres, a fibre of a degree-two map containing at most
one ramification point. The branch divisor of $\pi$ is a relative effective Cartier divisor of
degree two on the smooth model of $C$ whose generic fibre is $B$ and whose special fibre is the branch
divisor of $\pi_p$; so the two points of $B$ reduce to distinct points of $C_p$. For $B=\{(a:\pm y':
b)\}$ with $\gcd(a,b)=1$ the two reductions are $(\bar a:\pm\bar y':\bar b)$, which coincide exactly
when $p\mid 2y'$. Hence $p\nmid y'$ for every $p\notin S$.

Now $y'^2=G(a,b)$ for the binary sextic form $G(a,b)=b^6g(a/b)$, so $G(a,b)$ is an $S$-unit. Over
$\Z$, $G=(8a+b)\,Q_5(a,b)$ with
\[
  Q_5(a,b)=8a^5+a^4b+90a^3b^2-239a^2b^3+128ab^4+16b^5
\]
irreducible, and the resultant of the two factors is $2^7\cdot7\cdot11\cdot13$; so any prime
dividing both lies in $S$, and both factors are $S$-units. Thus $(a,b)$ is a coprime solution of
the Thue--Mahler equation
\[
  Q_5(a,b)=\pm\,2^{z_1}7^{z_2}11^{z_3}13^{z_4}.
\]
Such equations have finitely many solutions, and the algorithm of Gherga and Siksek \cite{gs} finds
them all: here it returns $21$ coprime pairs $(a,b)$ up to sign (\texttt{stepH-thuemahler.sols}),
among them $(1,0)$, $(0,1)$, $(3,2)$ and $(1,3)$. For each, $G(a,b)$ is a square with $S$-unit
root exactly for the seven pairs listed (\texttt{stepH-thuemahler.m}); the pair $(1,-8)$, where
$G=0$, is the Weierstrass point $W$ and not a fibre of two points, and the pairs such as $(3,1)$,
where $G=260^2$, have $8a+b=25$ and fail.
\end{proof}

Along each of the seven fibres the $32$ twists are now compared with $\#Y(\F_p)$ at $p=3,5,17,19,
23,29,31,37,41,43,47$: exactly one pair survives, the fibre at infinity with $c=2$
(\texttt{stepH-thuemahler.m}; \texttt{stepH-enumerate.m} runs the same comparison over the fibres
of all rational points of height at most $10^6$, which are the same seven and three more, with the
same outcome). And $y^2=q_1(t^2/2)$ is $Y_2$ after $y\mapsto4y$, while $v^2=2u\,q_1(u)$ is $C'$,
and $Y_2$ has the counts of $Y$ at all $31$ good primes below $150$ (\texttt{stepH-verify.m}). This
proves (3). \qed

\subsection*{Hyperellipticity, exactly}
By (2), $Y$ is hyperelliptic exactly when $P_{11}+P_{91}$ is a fibre of the hyperelliptic map of
$C$, that is, when it lies in the canonical class; and on a curve of genus two two \emph{distinct}
effective divisors of degree two that are linearly equivalent are both canonical, since
$h^0(D)\ge2$ forces $D\in|K_C|$. This turns the question into one about a divisor class on the
star, where Borcherds products live and where the counting argument of the note on $X_0^{46}(31)^*$
applies: a Borcherds form has a divisor class of finite order \cite{borcherds}, and a class of
finite order in $\Jac(C)(\Q)$ is trivial when that group has no torsion.

\begin{proposition}\label{prop:count}
$\Jac(C)(\Q)$ is torsion-free. Consequently, if for some $n\ge1$ there is a Borcherds form on
$X_0^{2002}(1)^*$ whose divisor is $n\bigl(Z(-11)+Z(-91)-Z(d)-Z(d')\bigr)$, where $Z(d)$ and
$Z(d')$ are rational CM points of $C$ distinct from $P_{11}$, $P_{91}$ and from each other, then
$P_{11}+P_{91}$ and $Z(d)+Z(d')$ are both fibres of the hyperelliptic map of $C$, $Y$ is
hyperelliptic, and $Y\cong Y_2$ over $\Q$.
\end{proposition}

\begin{proof}
The torsion subgroup of $\Jac(C)(\Q)$ is computed exactly from the model (\texttt{stepH-verify.m}):
it is trivial. A Borcherds form of weight zero on $X_0^{2002}(1)$ with that divisor is a rational
function whose divisor is $G$-invariant; its $|G|$-th power descends to $C$, where the Heegner
divisors $Z(-11)$, $Z(-91)$, $Z(d)$, $Z(d')$, each a single $G$-orbit, become the points
$P_{11},P_{91},Z(d),Z(d')$ with multiplicities given by the ramification of $X_0^{2002}(1)\to C$ at
each, which are positive integers. So some positive multiple $m\bigl(P_{11}+P_{91}-Z(d)-Z(d')\bigr)$
is the divisor of a function on $C$ over $\Q$ --- the class is Galois-stable, the form's divisor
being defined over $\Q$ --- hence the class of $P_{11}+P_{91}-Z(d)-Z(d')$ is torsion in
$\Jac(C)(\Q)$, hence zero. The two effective divisors are distinct and linearly equivalent, so both
are canonical: $P_{91}=\bar\iota P_{11}$, and Theorem~\ref{thm:genus4}(2),(3) finish.
\end{proof}

The CM points of the star, located on the model by the Shimizu lift (and all rational), are:
$P_{-67}$ over $x=1$, $P_{-99}$ over $x=\tfrac13$, $P_{-148}$ and $P_{-232}$ both over $x=0$,
$P_{-163}$ and $P_{-312}$ both over $x=3$, $P_{-280}$ over $x=-1$; each of $Z(-148)$, $Z(-232)$,
$Z(-163)$, $Z(-312)$ is a single point by Ogg's counts (sixteen points of the base in one free
orbit). So the natural candidates are $Z(-11)+Z(-91)-Z(-148)-Z(-232)$ and $Z(-11)+Z(-91)-Z(-163)
-Z(-312)$, each a difference of two numerical fibres, and the search for a Borcherds form with either
divisor, by the prescribed-divisor search used for Guo--Yang's Example 43 and for $X_0^{46}(31)^*$,
is running at the time of writing. % BORCHERDS-PENDING: lovelace ~/g2borch/borch/b2002_bd{2,4}.log

\begin{remark}\label{rem:gap}
The numerical evidence for hyperellipticity is of two kinds. Noether's count on the four
Shimizu-lifted differentials of $Y$ finds the three quadratic relations of a rational normal
quartic. And on $Y$ itself, where the forms do not vanish at the fixed points of $\bar w$ (the
covering $X_0^{2002}(1)\to Y$ is unramified there), the canonical images of $\tilde P_{11}$ and
$\tilde P_{91}$ are proportional to $1.3\cdot10^{-16}$ at $80$ digits, while unrelated pairs of CM
points give $0.5$ to $0.9$: so $\tilde P_{91}=\iota\tilde P_{11}$ numerically, which is
hyperellipticity by Lemma~\ref{lem:normalform}. Neither can become a proof, because hyperellipticity
is the vanishing of a class, which numerics can only bound; Proposition~\ref{prop:count} is the
exact replacement. Note that Proposition~\ref{prop:fibre} assumed $B$ to be a fibre: for a
non-hyperelliptic $Y$ its two rational branch points would only be constrained not to coincide
modulo any prime outside $S$, an $S$-integrality condition on $C\times C$ minus the diagonal with no
effective finiteness, which is why the exact route to hyperellipticity goes through the divisor
class and not through the branch points.
\end{remark}

\begin{remark}
The same shape recurs. Each of the open quotients $X_0^{2030}(1)/H$, $X_0^{2418}(1)/H$,
$X_0^{3390}(1)/H$, $X_0^{6}(341)/H$ and $X_0^{6}(403)/H$ (with $H$ of index two in the
Atkin--Lehner group) has genus four over a star of genus two, so once the star is proved by the
method of Section~2 the present section applies verbatim: the branch locus is read off Ogg's
formula, hyperellipticity is the position of one rational divisor, and the model, if any, is
$y^2=q_1(t^2/c)$ with $c$ pinned by counts.
\end{remark}

\section{The computations}\label{sec:computations}

Each line is a script in \texttt{tools/proof-2002/}, and each prints the number quoted.
\begin{itemize}
\item \texttt{stepH-branch.m}: the fixed points of every $w_m$, $m\notin H$, on $X_0^{2002}(1)$ by
  discriminant, and the genera of $Y$ and of the star.
\item \texttt{stepH-enumerate.m}: the rational points of $C$ of height at most $10^6$, the
  $S$-unit test on each fibre, and for each fibre and each of the $32$ twists the counts of
  $C'_c$ against those forced by the trace formula.
\item \texttt{stepH-verify.m}: $Y_2$ against the trace formula at the $31$ good primes below $150$,
  its two quotients against $C$ and $C'$, the quintic model $q_1$, the torsion of $\Jac(C)(\Q)$, and
  the conductor of $C'$.
\item \texttt{stepH-thuemahler.m}, with \texttt{stepH-thuemahler.sols}: the $21$ solutions of the
  Thue--Mahler equation as returned by the solver of \cite{gs} (run outside this repository, in
  $237$ s), each re-verified; the seven fibres they give; and the twist-and-count test along each.
\item \texttt{step0-conductor-euler.m}: the odd part of the conductor, and the Euler factors of
  $\Jac(C)$ against $1-t_pT+(n_p+2p)T^2-t_ppT^3+p^2T^4$ at every good prime below $120$. ($t_p$
  and $n_p$ are read off the characteristic polynomial of $T_p$ on the four-dimensional space of
  modular symbols of $f$, which is $(x^2-t_px+n_p)^2$; its square root is taken as a polynomial,
  since a factorisation returns a linear factor at the primes where $a_p$ is rational.)
\item \texttt{step0-conductor-at-2.m}: the exponent at $2$, by \texttt{EvenConductorExponent\_Genus2}
  with \texttt{Proof := true} \cite{dd}.
\item \texttt{step2-endomorphisms.m}: \texttt{CertifiedEndomorphismAlgebra} \cite{cmsv} over $\Q$,
  with its output, the certificate, kept beside it.
\item \texttt{stepA-orbit.m}: the five orbits at level $2002$ and their Euler factors.
\item \texttt{stepB-reducibility.m}: $\gcd_p \Norm(a_p-1-p)=1$ over the $24$ primes $p\equiv1
  \pmod{2002}$ below $200000$.
\item \texttt{stepB-congruence.m}: $\gcd_p (a_p-\sigma a_p)^2=5$ over the good primes below $400$.
\item \texttt{step5-ramified.m}: the group structure of $\Jac(C)(\F_p)$ for $p<200$, and the primes
  of odd $5$-rank.
\end{itemize}

\begin{thebibliography}{9}
\bibitem{borcherds} R.~E. Borcherds, \emph{The Gross--Kohnen--Zagier theorem in higher dimensions},
Duke Math. J. \textbf{97} (1999), 219--233; correction, Duke Math. J. \textbf{105} (2000), 183--184.
\bibitem{bppt} A.~Brumer, A.~Pacetti, C.~Poor, G.~Tornar\'ia, J.~Voight and D.~S. Yuen, \emph{On
the paramodularity of typical abelian surfaces}, Algebra Number Theory \textbf{13} (2019),
1145--1195.
\bibitem{carayol} H.~Carayol, \emph{Sur les repr\'esentations $\ell$-adiques associ\'ees aux formes
modulaires de Hilbert}, Ann. Sci. \'Ecole Norm. Sup. \textbf{19} (1986), 409--468.
\bibitem{carayol89} H.~Carayol, \emph{Sur les repr\'esentations galoisiennes modulo $\ell$
attach\'ees aux formes modulaires}, Duke Math. J. \textbf{59} (1989), 785--801.
\bibitem{cmsv} E.~Costa, N.~Mascot, J.~Sijsling and J.~Voight, \emph{Rigorous computation of the
endomorphism ring of a Jacobian}, Math. Comp. \textbf{88} (2019), 1303--1339.
\bibitem{dd} T.~Dokchitser and C.~Doris, \emph{3-torsion and conductor of genus 2 curves}, Math.
Comp. \textbf{88} (2019), 1913--1927.
\bibitem{edixhoven} B.~Edixhoven, \emph{The weight in Serre's conjectures on modular forms},
Invent. Math. \textbf{109} (1992), 563--594.
\bibitem{gs} A.~Gherga and S.~Siksek, \emph{Efficient resolution of Thue--Mahler equations},
preprint, arXiv:2207.14492 (2022).
\bibitem{ggr} J.~Gonz\'alez, J.~Gu\`ardia and V.~Rotger, \emph{Abelian surfaces of $\mathrm{GL}_2$-type
as Jacobians of curves}, Acta Arith. \textbf{116} (2005), 263--287.
\bibitem{kw} C.~Khare and J.-P. Wintenberger, \emph{Serre's modularity conjecture, I, II}, Invent.
Math. \textbf{178} (2009), 485--504, 505--586.
\bibitem{lmfdb} The LMFDB Collaboration, \emph{The $L$-functions and modular forms database},
newform \texttt{2002.2.a.f}, \url{https://www.lmfdb.org/ModularForm/GL2/Q/holomorphic/2002/2/a/f/}.
\bibitem{ogg} A.~P. Ogg, \emph{Real points on Shimura curves}, in: Arithmetic and geometry, Vol.~I,
Progr. Math. \textbf{35}, Birkh\"auser, 1983, 277--307.
\bibitem{ribet80} K.~A. Ribet, \emph{Sur les vari\'et\'es ab\'eliennes \`a multiplications
r\'eelles}, C. R. Acad. Sci. Paris \textbf{291} (1980), 121--123; and \emph{On modular
representations of $\mathrm{Gal}(\overline{\Q}/\Q)$ arising from modular forms}, Invent. Math.
\textbf{100} (1990), 431--476, \S4.
\bibitem{ribet92} K.~A. Ribet, \emph{Abelian varieties over $\Q$ and modular forms}, in: Algebra and
topology 1992 (Taej\u{o}n), Korea Adv. Inst. Sci. Tech., 1992, 53--79.
\bibitem{serre-lauter} J.-P. Serre, appendix to K.~Lauter, \emph{Geometric methods for improving
the upper bounds on the number of rational points on algebraic curves over finite fields}, J.
Algebraic Geom. \textbf{10} (2001), 19--36.
\end{thebibliography}
\end{document}
